\documentclass[11pt]{article}
\usepackage[margin=0.85in]{geometry}
\usepackage[T1]{fontenc}
\usepackage{lmodern}
\usepackage{amsmath,amssymb,bm,booktabs,microtype}
\usepackage[colorlinks=true,linkcolor=blue,citecolor=blue,urlcolor=blue]{hyperref}
\hypersetup{pdftitle={Retained adjoint-fermion binding with explicit hbar},pdfauthor={Research notebook},pdfsubject={Version 0.13: nine dimensions, winding labels, exceptional groups, strings as points}}
\setlength{\parskip}{3pt}
\setlength{\emergencystretch}{2em}
\title{Retained adjoint-fermion binding with explicit $\hbar$:\\
recoil, quantum corrections and the whole-state current}
\author{Research-notebook draft, version 0.13}
\date{11 November 2026}
\begin{document}
\maketitle
\begin{abstract}
We study two fermionic endpoints bound to a retained adjoint fermion by regulated contact interactions, with no direct quark--antiquark force. In a fixed-number, global-color model, dimer binding suffices to prove a trimer below breakup. With physical units and explicit $\hbar$, we distinguish recoil perturbation theory, the semiclassical expansion and contact iterations. The leading recoil energy correction is second order and attractive. For a Gaussian contact in the deep-well regime, we derive the energy and whole-state charge radius through order $\hbar^2$, with exact mass-ratio dependence at each displayed order. The leading anharmonic radius correction is positive. Contact binding does not by itself make this current pointlike. A formal $\pi/K$ scalar pair supplies the charged Dirac muon multiplet count, while exact physical $B-L$ imposes a further conditional charge constraint. We compare a distinct 1983 composite-fermion mass-protection mechanism and retain endpoint and Goldstino proposals as construction hypotheses. A common string realization remains a separate task.
\end{abstract}

\paragraph{Research scope and provenance.}
The broader target remains bosonic strings with $q/\bar q$ or $q/q$ endpoints for mesons or diquarks, and fermionic strings with $q/\bar q$ or $\bar q/\bar q$ endpoints for leptons or quarks. The main phenomenological question is approximately pointlike scattering of the whole fermionic state. The user now selects the retained adjoint contact model as the most interesting concrete object and requests explicit $\hbar$ in all calculations.

This focused revision develops that selection. It does not identify a contact trimer with a string by definition. Version 0.2, including the warped-current comparison and the scalar-completion study, remains preserved at repository commit \texttt{6b73d29}; version 0.1 is preserved at \texttt{333e537}. The uploaded originals remain unchanged. Earlier natural-unit calculations are historical comparisons, not the unit convention of this paper. The new radius calculation uses standard perturbation theory; no claim of literature priority follows from its application here. The later hypothesis section does not change the fermionic-endpoint Hamiltonian used for this calculation.

\section{A retained adjoint binds both endpoints}
\label{sec:model}
We take distinguishable nonrelativistic endpoint species $q,\bar q$ with physical masses $m,m$, and a retained adjoint fermion $\lambda$ with physical mass $M$. Constituent numbers are conserved and color is global $SU(3)$. These assumptions define the model, not a relativistic QCD completion. The carrier is present in the bound state; it is not a virtual fermion exchanged between unchanged endpoints.

Throughout, $\hbar$ is explicit, $c_{\rm light}$ denotes the speed of light, positions have units of length and momenta have units of momentum:
\begin{equation}
 \mathbf p=-i\hbar\nabla,\qquad
 i\hbar\partial_t\Psi=H\Psi,\qquad [G]={\rm energy}\,{\rm length}^3 .
\end{equation}
The color matrices are dimensionless and obey
$\mathrm{tr}(T^aT^b)=\delta^{ab}/2$. The unique singlet in
$\mathbf3\otimes\bar{\mathbf3}\otimes\mathbf8$ has
\begin{equation}
 T_q\cdot T_{\bar q}=\frac16,\qquad
 T_q\cdot T_\lambda=T_{\bar q}\cdot T_\lambda=-\frac32.
\end{equation}
There is no direct $q\bar q$ potential. With
$\mathbf x=\mathbf r_q-\mathbf r_\lambda$ and
$\mathbf y=\mathbf r_{\bar q}-\mathbf r_\lambda$, set
\begin{align}
 f_R(r)&=\frac{e^{-r^2/R^2}}{\pi^{3/2}R^3},&
 V&=-\frac{3G}{2}[f_R(x)+f_R(y)],\qquad G>0,\\
 \mu&=\frac{mM}{m+M},&
 H_{\rm rel}&=\frac{p_x^2+p_y^2}{2\mu}
       +\frac{\mathbf p_x\cdot\mathbf p_y}{M}+V .
 \label{eq:hamiltonian}
\end{align}
Equivalently the kinetic operator is
$-\hbar^2(\nabla_x^2+\nabla_y^2)/(2\mu)
-\hbar^2\nabla_x\cdot\nabla_y/M$.
The two contacts are number-conserving four-fermion density interactions. No propagating bosonic binding field is included in this Hamiltonian. Their ultraviolet origin remains unspecified. $R>0$ is the specified finite regulator length; we do not silently take an unrenormalized three-dimensional delta limit.

The spin-independent model allows the two spin-$1/2$ endpoints to form a spin singlet, with total spin $1/2$ supplied by $\lambda$. Each constituent has $s^2=3\hbar^2/4$; the endpoints here are not scalar particles. Spin-dependent interactions, allowed annihilation channels and gauged color would change the problem and must be supplied explicitly.

\subsection{Binding and the correct breakup threshold}
For $h_\mu=p^2/(2\mu)-3Gf_R/2$, a normalized Gaussian orbital of width $a$ gives
\begin{equation}
 E_2^{\rm trial}(a)=\frac{3\hbar^2}{4\mu a^2}
   -\frac{3G}{2\pi^{3/2}(R^2+a^2)^{3/2}}.
\end{equation}
At $a^2=2R^2$ this proves the sufficient condition
\begin{equation}
 \boxed{G>\frac{3\sqrt3\,\pi^{3/2}}{4}\frac{\hbar^2R}{\mu}
 \quad\Longrightarrow\quad \epsilon_2<0 .}
 \label{eq:dimer-sufficient}
\end{equation}
Here $\epsilon_2$ is the exact dimer ground energy, not its variational estimate. If $u$ is its real normalized orbital, the product $u(x)u(y)$ has
$\langle\mathbf p_x\cdot\mathbf p_y\rangle=0$ and energy $2\epsilon_2$. Both dimer-plus-spectator thresholds are $\epsilon_2$, so
\begin{equation}
 \boxed{E_3\leq2\epsilon_2<\epsilon_2<0.}
 \label{eq:binding}
\end{equation}
The bounded short-range potentials and variational principle establish a trimer below breakup. Removing $\lambda$ removes both interactions. This is a binding result with no pre-existing $q\bar q$ bound state.

Nonrelativistic control requires characteristic momenta much smaller than $mc_{\rm light}$ and $Mc_{\rm light}$, and binding energies small compared with the corresponding rest energies. A useful range choice is $R\gg\hbar/(mc_{\rm light}),\hbar/(Mc_{\rm light})$. The model does not yet generate a light relativistic lepton, a confining flux tube, or protection against annihilation in a larger theory.

\section{Three different expansions}
Define a depth, a dimensionless contact strength, and a recoil parameter:
\begin{equation}
 V_0=\frac{3G}{2\pi^{3/2}R^3},\qquad
 \gamma=\frac{2\mu V_0R^2}{\hbar^2}
       =\frac{3\mu G}{\pi^{3/2}\hbar^2R},\qquad
 \eta=\frac{\mu}{M}=\frac{m}{m+M}.
 \label{eq:parameters}
\end{equation}
The energy unit $E_R=\hbar^2/(2\mu R^2)$ and coordinates
$\mathbf X=\mathbf x/R$, $\mathbf Y=\mathbf y/R$ give
\begin{equation}
 \frac{H_{\rm rel}}{E_R}
 =-\nabla_X^2-\nabla_Y^2-2\eta\nabla_X\cdot\nabla_Y
   -\gamma(e^{-X^2}+e^{-Y^2}).
\end{equation}
This rescaling displays $\hbar$ in $E_R$ and $\gamma$; it does not set it to one. Equation~\eqref{eq:dimer-sufficient} becomes $\gamma>9\sqrt3/4$.

Three expansions must be kept distinct: a recoil series in $\eta$ about two bound dimers; a deep-well semiclassical series in
$\varepsilon=\gamma^{-1/2}\propto\hbar$ at fixed physical potential and masses; and a Born series in the contact interaction about free propagation. They have different domains. In particular, the existence of a shallow bound-state pole generally requires resumming the Born series.

\section{Recoil perturbation theory about the bound product}
\label{sec:recoil}
At fixed $\mu,G,R$ write
\begin{equation}
 H(\eta)=H_0+\eta W,\qquad
 H_0=h_\mu(x)+h_\mu(y),\qquad
 W=\frac{\mathbf p_x\cdot\mathbf p_y}{\mu}.
\end{equation}
The unperturbed orbital ground state is $|0\rangle=|u,u\rangle$, with $E_0=2\epsilon_2$. Spin degeneracy is inert in this spin-independent Hamiltonian. Reversing $\mathbf x$ alone maps $H(\eta)$ unitarily to $H(-\eta)$. Thus the isolated orbital ground-energy branch is even near $\eta=0$:
\begin{equation}
 E_3(\eta)=2\epsilon_2-\eta^2\mathcal C_2+O(\eta^4),\qquad
 \mathcal C_2=
 \frac{1}{\mu^2}\sum_{a,b}
 \frac{|\langle a,b|\mathbf p_x\cdot\mathbf p_y|u,u\rangle|^2}
      {\Delta_a+\Delta_b}>0.
 \label{eq:recoil-second}
\end{equation}
Here $\Delta_a=\epsilon_a-\epsilon_2$; the spectral sum includes integrals over continuum states. Each momentum takes the dimer $S$ wave into an $l=1$ state, and their scalar contraction selects total orbital angular momentum zero. The first-order shift vanishes, while the second-order shift is attractive. The intermediate states need not be bound $P$ waves. The coefficient has units of energy, and its momenta contain the explicit $\hbar$ above.

This is a correction from the motion of the same retained carrier. It does not insert a direct pair potential between the endpoints. The product trial state already proves binding; the correlated wavefunction lowers its energy further. Recoil correlations are an actual dynamical effect even though their first-order energy expectation vanishes.

For subsequent orders, let $\mathcal Q=1-|0\rangle\langle0|$ and
$\mathcal R_0=\mathcal Q(E_0-H_0)^{-1}\mathcal Q$.
With intermediate normalization and
$|\psi\rangle=\sum_{n\geq0}\eta^n|\psi_n\rangle$, one has
\begin{align}
 E_n&=\langle0|W|\psi_{n-1}\rangle,\\
 |\psi_n\rangle&=\mathcal R_0\left[
 W|\psi_{n-1}\rangle-\sum_{j=1}^{n-1}E_j|\psi_{n-j}\rangle
 \right],\qquad n\geq1 .
 \label{eq:recursion}
\end{align}
Odd energy coefficients vanish. This is an all-order prescription, not an evaluation of every coefficient or a claim of convergence up to $\eta=1$. A small dimer gap can limit its useful domain.

The fixed-$\mu$ qualification matters: changing $M$ at fixed endpoint mass $m$ also changes $\mu$. The dimer contribution $2\epsilon_2(\mu)$ then changes, and the total physical heavy-mass expansion need not contain only even powers of $1/M$.

\section{An analytic deep-well calculation through order \texorpdfstring{$\hbar^2$}{hbar squared}}
\label{sec:deepwell}
For the specified Gaussian contact take $\gamma\gg1$. Define
\begin{equation}
 \omega^2=\frac{2V_0}{\mu R^2},\qquad
 \varepsilon=\frac{\hbar}{\mu\omega R^2}=\gamma^{-1/2}\ll1,\qquad
 S_\eta=\sqrt{1+\eta}+\sqrt{1-\eta}.
\end{equation}
Hold $0\leq\eta<1$ fixed in this expansion. Deep binding relative to $E_R$ can coexist with nonrelativistic control if $V_0$ remains small compared with the constituent rest energies and $R$ is sufficiently large.

In $\mathbf u_\pm=(\mathbf x\pm\mathbf y)/\sqrt2$, the harmonic Hamiltonian is
\begin{equation}
 H_{\rm harm}=-2V_0+
 \sum_{\sigma=\pm}\left[
 \frac{(1+\sigma\eta)p_\sigma^2}{2\mu}
 +\frac{\mu\omega^2u_\sigma^2}{2}
 \right].
\end{equation}
Its frequencies are $\omega\sqrt{1\pm\eta}$, and hence
\begin{equation}
 E_{3,{\rm harm}}=-2V_0+\frac{3\hbar\omega}{2}S_\eta .
 \label{eq:harmonic}
\end{equation}
This resums recoil exactly within the harmonic approximation. Relative to twice the harmonic dimer energy,
\begin{equation}
 E_{3,{\rm harm}}-2\epsilon_{2,{\rm harm}}
 =-\frac{3\hbar\omega}{8}\eta^2
  -\frac{15\hbar\omega}{128}\eta^4+O(\hbar\omega\eta^6).
\end{equation}

The first anharmonic term is
$V_4=-V_0(x^4+y^4)/(2R^4)$, where $x^4=(\mathbf x^2)^2$.
The variance of each Cartesian component of $\mathbf x$ or $\mathbf y$ in the harmonic ground state is
$s=\hbar S_\eta/(4\mu\omega)$. Since $\langle x^4\rangle=15s^2$, first-order perturbation theory in $V_4$ gives
\begin{equation}
 \boxed{
 E_3=-2V_0+\frac{3\hbar\omega}{2}S_\eta
 -\frac{15\hbar^2}{32\mu R^2}S_\eta^2
 +O(V_0\varepsilon^3).}
 \label{eq:hbar-result}
\end{equation}
Thus the classical well depth, harmonic zero-point energy and first anharmonic correction are explicitly orders $\hbar^0,\hbar^1,\hbar^2$ at fixed $V_0,R,\mu,\eta$. Sextic first-order perturbation and quartic second-order perturbation both first enter at order $\hbar^3$; retaining only one would be incomplete.

The corresponding dimer expansion is
\begin{equation}
 \epsilon_2=-V_0+\frac{3\hbar\omega}{2}
       -\frac{15\hbar^2}{16\mu R^2}
       +O(V_0\varepsilon^3).
\end{equation}
In particular, the coefficient of $\eta^2$ in the combined expansion is
\begin{equation}
 E_3-2\epsilon_2=
 \left[-\frac{3\hbar\omega}{8}
       +\frac{15\hbar^2}{32\mu R^2}
       +O(\hbar\omega\varepsilon^2)\right]\eta^2
       +O(\hbar\omega\eta^4).
\end{equation}
The attractive curvature agrees with Eq.~\eqref{eq:recoil-second} in the controlled regime $\varepsilon\ll1$. These formulas concern the chosen finite Gaussian regulator; they are not universal coefficients of every short-range force.

\section{The whole-state current with physical momentum transfer}
\label{sec:current}
Take an electrically neutral $\lambda$ and endpoint charges $e_1,e_2$ with $e_1+e_2\ne0$. Equal endpoint masses and symmetric dynamics make their one-body elastic form factors equal, even when their charges differ. A charged state may require different endpoint flavors; the binding proof did not require a nonzero electric charge.

The leading constituent current measures
\begin{equation}
 \mathbf r_q-\mathbf R_{\rm cm}
 =\frac{(m+M)\mathbf x-m\mathbf y}{2m+M}
 =\frac{\mathbf x-\eta\mathbf y}{1+\eta},\qquad
 F(\mathbf Q)=
 \left\langle e^{i\mathbf Q\cdot(\mathbf r_q-\mathbf R_{\rm cm})/\hbar}
 \right\rangle .
\end{equation}
Here $\mathbf Q$ is physical momentum transfer, not wave number. With
$r_E^2=\langle|\mathbf r_q-\mathbf R_{\rm cm}|^2\rangle$,
\begin{equation}
 F(Q)=1-\frac{Q^2r_E^2}{6\hbar^2}
       +O\!\left(\frac{Q^4\langle r^4\rangle}{\hbar^4}\right).
 \label{eq:smallq}
\end{equation}
For a neutral total state the normalized electric form factor is unavailable; other currents or multipoles must be specified.

The harmonic correlated state gives
\begin{align}
 \langle x^2\rangle_{\rm harm}
 &=\frac{3\hbar}{4\mu\omega}S_\eta,&
 \langle\mathbf x\cdot\mathbf y\rangle_{\rm harm}
 &=\frac{3\hbar}{4\mu\omega}D_\eta,\\
 D_\eta&=\sqrt{1+\eta}-\sqrt{1-\eta},&
 r_{E,{\rm harm}}^2
 &=\frac{3\hbar[(1+\eta^2)S_\eta-2\eta D_\eta]}
         {4\mu\omega(1+\eta)^2}.
\end{align}
The positive cross-correlation is generated by carrier recoil. It reduces the mean-square radius relative to a distribution with the same marginal variance but zero cross-correlation. The total change also includes the change in that marginal variance and in the center-of-mass weights.

The Gaussian harmonic result is
\begin{equation}
 F_{\rm harm}(Q)=\exp\!\left[-\frac{Q^2r_{E,{\rm harm}}^2}{6\hbar^2}\right].
 \label{eq:gaussian-current}
\end{equation}
For a single dressed current insertion, after matching the point-particle kinematics and assuming no independent magnetic or transition form factors, the Born rate ratio is $|F|^2$. The illustrative condition $|F|^2\geq1-\delta$ gives
\begin{equation}
 r_{E,{\rm harm}}^2\leq
 -\frac{3\hbar^2}{Q^2}\log(1-\delta),\qquad 0<\delta<1 .
\end{equation}
Contact binding alone therefore does not remove size from the current. In particular, a mean-square radius proportional to $\hbar$ is insufficient by itself for pointlike scattering at fixed physical $Q$: the resolution length is $\hbar/Q$. The harmonic approximation to the current is controlled for fixed $Q\sqrt{r_E^2}/\hbar$ in the deep-well expansion, especially near the pointlike window. Its exponentially small large-$Q$ tail is not a uniform prediction of the finite Gaussian well as $\hbar\to0$.

\subsection{Leading anharmonic radius}
\label{sec:anharmonic-radius}
The first correction to the current's slope requires the changed bound state.
It can be obtained without a spectral sum by adding independent quadratic
sources $j_\pm u_\pm^2/2$, writing $k_\pm=\mu\omega^2+j_\pm$, and
holding the physical quartic coefficient $-V_0/(2R^4)$ fixed. Define the component variances
\begin{equation}
 v_\pm=\frac{\hbar}{2}\sqrt{\frac{1\pm\eta}{\mu k_\pm}},
 \qquad s=\frac{v_++v_-}{2}.
\end{equation}
The harmonic energy depending on these sources is
$(3\hbar/2)\sum_\pm\sqrt{(1\pm\eta)k_\pm/\mu}$, and
$E_4=-15V_0s^2/R^4$. Hellmann--Feynman differentiation gives
\begin{equation}
 \langle u_\pm^2\rangle
 =2\frac{\partial(E_{\rm harm}+E_4)}{\partial k_\pm}
 =3v_\pm+\frac{15V_0s v_\pm}{k_\pm R^4}
 +O(R^2\varepsilon^3).
\end{equation}
Evaluate at $k_+=k_-=\mu\omega^2=2V_0/R^2$.
In the full finite well these derivatives may be taken from nonnegative
added sources, avoiding an unbounded negative quadratic potential at infinity.
Exchange symmetry makes $\langle\mathbf u_+\cdot\mathbf u_-\rangle=0$.
Both mode variances, and hence both $\langle x^2\rangle$ and
$\langle\mathbf x\cdot\mathbf y\rangle$, acquire the same relative
correction $5S_\eta\varepsilon/8$. Thus
\begin{equation}
 \boxed{r_E^2=r_{E,{\rm harm}}^2
 \left(1+\frac{5S_\eta}{8}\frac{\hbar}{\mu\omega R^2}\right)
 +O(R^2\varepsilon^3).}
 \label{eq:radius-correction}
\end{equation}
The leading radius is order $\hbar$ and this correction is order $\hbar^2$
at fixed physical well parameters and mass ratio. The coefficient is positive:
the finite Gaussian well is softer than its quadratic approximation and
increases the size. It does not supply an additional pointlike mechanism.
In the independent-dimer limit $\eta=0$,
\begin{equation}
 r_E^2=\frac{3\hbar}{2\mu\omega}
 \left(1+\frac{5\hbar}{4\mu\omega R^2}\right)
 +O(R^2\varepsilon^3),
\end{equation}
which agrees with the single three-dimensional Gaussian well. The coefficients
are at fixed $0\leq\eta<1$; no uniform limit at $\eta=1$ is asserted.
Equation~\eqref{eq:radius-correction} improves the small-$Q$ expansion
\eqref{eq:smallq}, under the same charge assumptions. Substituting this radius
into a Gaussian exponential would not derive the anharmonic finite-$Q$
form factor: its higher moments change as well.

\section{Contact iterations are not the same as the semiclassical series}
Choose momentum normalization
$\langle\mathbf r|\mathbf p\rangle=e^{i\mathbf p\cdot\mathbf r/\hbar}$, with completeness measure $d^3p/(2\pi\hbar)^3$. The dimer potential kernel is
\begin{equation}
 V(\mathbf p',\mathbf p)
 =-\frac{3G}{2}
   \exp\!\left[-\frac{R^2|\mathbf p'-\mathbf p|^2}{4\hbar^2}\right].
\end{equation}
Its scattering operator obeys
\begin{equation}
 T(\mathbf p',\mathbf p;E)=V(\mathbf p',\mathbf p)
 +\int\frac{d^3k}{(2\pi\hbar)^3}
 \frac{V(\mathbf p',\mathbf k)\,T(\mathbf k,\mathbf p;E)}
      {E-k^2/(2\mu)+i0}.
 \label{eq:ls}
\end{equation}
Every factor has physical units: $V$ and $T$ here have units of energy times volume. Repeated propagation in this equation is organized by the contact strength and the energy denominator, not by assigning an extra positive power of $\hbar$ to every iteration. At fixed physical $G,R,\mu$, $\gamma\propto\hbar^{-2}$; taking $\hbar$ small makes this interaction strong in dimensionless Schr\"odinger units.

A bound-state pole is a nonperturbative feature of the full resolvent. Quantum-field-theory loop counting instead follows from a specified action in $\exp(iS/\hbar)$ and a stated field/coupling normalization. The present fixed-number Hamiltonian has no vacuum pair-creation loops. A relativistic completion would have to specify them before their $\hbar$ orders could be calculated. We have not computed such a completion's radiative corrections.

\section{Result and next calculation}
The retained adjoint contact model binds both endpoints without a direct endpoint force. Keeping $\hbar$ explicit reveals three separate controls, attractive second-order recoil, and the deep-well energy and whole-state radius through $\hbar^2$. The leading anharmonic radius correction is positive. This crosses the current calculation's milestone without establishing the desired experimental pointlikeness.

The next quantitative target is the first anharmonic correction to the finite-momentum current for fixed $Q\sqrt{r_{E,{\rm harm}}^2}/\hbar=O(1)$, including its non-Gaussian shape. This will test a declared resolution interval rather than only its zero-momentum slope. A different near-threshold target would require resummed two- and three-body propagation. Neither needs an immediate three-generation model. A string or Seiberg-dual completion remains a possible bridge to the broader correspondence, with an explicit operator dictionary and physical currents.

The previous scalar-completion study produces spin-dependent contacts rather than the spin-independent interaction assumed here. For physical spins its projectors are
$P_0=1/4-\mathbf s_i\cdot\mathbf s_j/\hbar^2$ and $P_1=1-P_0$.
Its binding proof cannot be imported unchanged. This limitation does not alter the result for the Hamiltonian in Eq.~\eqref{eq:hamiltonian}.

\section{Endpoint statistics and a charged supercharge: construction hypotheses}
\label{sec:susy-dictionary}
The string interpretation leaves an additional choice open. In an RNS
construction the physical NS/R sectors determine the whole string's spacetime
statistics; Chan--Paton labels alone do not specify independently propagating
spin-$1/2$ endpoint particles. Physical endpoints require boundary fields and
an action. The contact model above chooses two fermionic endpoints, while the
user now also proposes two scalar endpoints for the fermionic string. Either
pair has even fermion parity and can give an odd state with an odd retained
carrier. Mixed scalar/fermion endpoints plus an odd carrier instead have even
total parity. The ordinary SQCD mesino is odd with mixed endpoints and no
additional odd carrier. These are distinct candidates to match, not
interchangeable descriptions already derived from a common theory.

Three charged Dirac leptons have the spin degrees of freedom of six complex
spin-zero fields. The six charged pseudoscalar mesons formed from two up-type
and three down-type hadronizing flavors provide this count. Direct antecedents
exist: Rivero's 2005 composite-boson proposal and 2009 \emph{sBootstrap}
discussion explicitly pursue the three-generation correspondence and the
pointlike-lepton difficulty~\cite{Rivero2005,Rivero2009}. These sources correct
the earlier limited historical search; neither the count nor this research
direction is claimed as new here. The user proposes the pairs
$\mu\leftrightarrow(\pi,K)$, $\tau\leftrightarrow(D,D_s)$ and
$e\leftrightarrow(B,B_c)$, with charge conjugates understood. A dynamical
supersymmetry representation, mass splittings and the top's absence from the
ordinary meson list remain to be explained.

\subsection{A common \texorpdfstring{$B-L$}{B minus L} shift}
For conventional dimensionless number charges, let $X=B-L$. A candidate
supercharge $\mathcal Q_+$ can have the same charge shift in both branches:
\begin{center}
\begin{tabular}{lccc}
\toprule
Transition & Initial $X$ & Final $X$ & Shift\\
\midrule
meson $\to$ antilepton & $0$ & $+1$ & $+1$\\
antidiquark $\to$ quark & $-2/3$ & $+1/3$ & $+1$\\
meson $\to$ lepton & $0$ & $-1$ & $-1$\\
diquark $\to$ antiquark & $+2/3$ & $-1/3$ & $-1$\\
\bottomrule
\end{tabular}
\end{center}
The first two transitions use $\mathcal Q_+$ and the last two its conjugate
$\mathcal Q_-$. The corresponding color representations agree. With suitable
flavors the electric charges can agree as well. Algebraically,
\begin{equation}
 [X,\mathcal Q_\pm]=\pm\mathcal Q_\pm,
 \qquad [\hbar X,\mathcal Q_\pm]=\pm\hbar\mathcal Q_\pm.
\end{equation}
The second expression makes explicit the convention for a dimensionful
generator; the first does not set $\hbar$ to one. If these are spacetime
supercharges, $X$ acts as an R symmetry. Their conjugate charges allow a
neutral translation generator in their anticommutator. This is an admissible
charge ledger, not yet a representation or interaction law.

The two positive shifts require respectively
$(\Delta B,\Delta L)=(0,-1)$ and $(1,0)$.
One fixed charge eigenoperator cannot implement both while carrying definite
charges under independently conserved $B$ and $L$. In a retained-carrier
realization, an exact nonzero continuous $X$ charge also requires distinguishable
carrier/anticarrier sectors; a self-conjugate Majorana carrier cannot carry it.

There is a further, useful constraint on a \emph{linear, exact} $N=1$
interpretation. Choose the convention that a left-chiral scalar of charge $x$
has Weyl partner charge $x-1$. A charged Dirac lepton uses left-handed Weyl
fields with $X=-1,+1$, so the corresponding scalars have $X=0,+2$.
Equivalently the two scalar particles with the lepton's electric charge have
$X=0,-2$. Two ordinary mesons, both $X=0$, do not fill this multiplet with
conserved physical $B-L$. Possible construction choices include an additional
scalar with two units of lepton number, breaking or reassigning the R charge,
or a nonlinear/emergent supersymmetry without these simple observed-state
multiplets. The constraint has these explicit premises; it is not an
exclusion of the broader string proposal.

\subsection{The pion--muon--Goldstino conjecture}
Volkov and Akulov proposed a neutrino as the Goldstone fermion of nonlinear
supersymmetry~\cite{VolkovAkulov}. Their supplied two-page paper does not
identify a pion--muon partner pair. We retain the additional conjecture
\begin{equation}
 \pi^-\longrightarrow\mu^-+\bar G,
 \qquad \bar G\ \hbox{tentatively identified with }\bar\nu_\mu,
\end{equation}
as a broken-supersymmetry scalar-partner transition. The assignments
$X(\mu^-)=-1$ and $X(\bar G)=+1$ conserve $B-L$; the emitted state is
electrically neutral. Nonlinear supersymmetry does not itself require a light
scalar partner, so the proposed pion identification is additional physical
content. Its next substantive test would specify the Goldstino coupling and
compare pion leptonic amplitudes, their helicity dependence and the
muon/electron channel ratio. No such amplitude or decay rate is calculated
in this revision.

\section{Partial compositeness and the sBootstrap comparison}
\label{sec:partialcomparison}
The supplied v1 sources of Agashe, Ricci and Sundrum~\cite{ARS2025}
and Agashe, Giudice, Rattazzi and Sundrum~\cite{AGRS2026} provide two
related approaches to hierarchical flavor. Both couple elementary matter
to operators of a new strong sector. Their operator dimensions generate
hierarchical mixing strengths from structureless ultraviolet values.
In the irrelevant-coupling, approximately fixed-point regime,
\begin{equation}
 \widehat\lambda_i(E)=\widehat\lambda_i(\Lambda_{\rm UV})
 \left(\frac{E}{\Lambda_{\rm UV}}\right)^{\gamma_i}.
\end{equation}
Here both scales are energies and the couplings are dimensionless;
this relation introduces no convention setting physical $\hbar$ to one.
Relevant mixings, including those needed for the large top Yukawa,
require the nonlinear RG treatment of the sources.
The nonsupersymmetric fermionic-operator exponent is
$\gamma=\Delta_f-5/2$, whereas the supersymmetric chiral-superfield
exponent is $\gamma=\Delta_s-2$. The chiral-multiplet relation
$\Delta_f=\Delta_s+1/2$ makes these the same scaling criterion.

\paragraph{Different protection mechanisms.}
The 2025 construction generates hierarchies in a flavor-anarchic UV
stage and then assumes a transition to strong dynamics with emergent
flavor and CP symmetries. Right-handed mixings flow towards universal
values, while left-handed mixings retain the flavor hierarchies.
Its MFV example admits compositeness around 10 TeV, with a less symmetric
variant around 5 TeV. A significant consequence is that right-handed
fermions of \emph{all} generations can have sizable compositeness:
light mass does not require both chiralities to be nearly elementary.
The Higgs is protected as a pseudo-Nambu--Goldstone boson.

The 2026 Super-Composite Higgs instead retains a flavor-anarchic strong
sector and combines it with supersymmetry. Composite Higgs superfields
emerge at a scale $m_*$, residual strong dynamics breaks SUSY at a lower
scale $m_S$, and superpartner masses occur at the soft scale.
These source symbols denote energy scales. Supersymmetry permits a
higher flavor/compositeness scale while protecting Higgs lightness;
holomorphy protects a vanishing supersymmetric $\mu$ term until SUSY
breaking. Gauge-mediation-like soft terms and restrictions on
supersymmetric operators help control flavor effects. Compositeness
also supplies the proposed $\mu$--$B_\mu$ mechanism.
Thus this is not simply the 2025 emergent-flavor construction with
SUSY added. The analyzed scales extend to multi-PeV compositeness,
and the authors discuss electron-EDM, gravitino and collider signatures
with model and cosmological assumptions. Residual tuning remains.

\paragraph{Comparison with the original sBootstrap.}
Rivero's proposal~\cite{Rivero2005,Rivero2009} identifies ordinary QCD
mesons and diquark/antidiquark channels as scalar-partner candidates.
This describes the earlier hadronic identification used in this
notebook. The subsequently supplied 2024 paper explicitly calls for
different high-energy bound states; Section~\ref{sec:so32completion}
corrects and refines that comparison.
Its six charged meson flavor states and conjugates arise from the five
hadronizing flavors $u,c,d,s,b$, with top excluded. The two new papers
instead posit a new strong sector and take three SM generations and
suitable composite-operator families as inputs. Their six up/down
operator channels, where present, are not a derivation of those six
QCD mesons. Substantial top compositeness under a new force also differs
from top's inability to form ordinary QCD mesons.

An ordinary spin-zero meson cannot be the fermionic operator in a
linear elementary-lepton mixing. A suitable odd operator, its quantum
numbers and its overlap must be supplied. The retained-adjoint operator
studied in this notebook is a candidate hypothesis, not an operator
identified or dynamically derived by either paper.
Moreover, in the 2026 conventional supersymmetric portal, elementary
and composite fields mix as whole superfields. In an unbroken,
canonically normalized linear realization, the scalar and fermion
components share that mixing rotation. It does not independently
choose a wholly QCD-composite scalar and an elementary fermion.
This statement concerns that realization, not every possible emergent
or nonlinear supersymmetry.

There is also a charge distinction: ordinary mesons have $B=L=0$,
whereas conventional sleptons share their lepton partners' lepton
number. The 2026 assumptions preserving $B$ and $L$ do not implement
the sBootstrap identification. The earlier conditional physical
$B-L$/R analysis remains relevant. Finally, an SQCD meson chiral
superfield has a scalar--scalar lowest component and a
scalar--fermion fermionic component. Identifying it with a QCD
pseudoscalar built from fermionic quarks requires physical operator
matching, not just the shared word ``meson.''

\paragraph{What transfers and what remains open.}
The useful ingredients are RG-generated mixing hierarchies, emergent
flavor protection, and separation of composite-state formation from
SUSY breaking. The 2026 hidden-SQCD/ISS toy illustrates staged dynamics
but explicitly has an unrealistic gauge-mediated spectrum with
suppressed gaugino masses. Neither paper supplies a fully derived
realistic microscopic completion or a QCD sBootstrap vacuum.
Its Goldstino belongs to a SUSY-breaking sector and becomes the
gravitino's longitudinal modes; it is not identified with a neutrino.

Neither work calculates the whole-state current of the proposed
fermionic color string or proves that it is pointlike while its
bosonic partner is extended. Small mixing or a high new-physics scale
can suppress observable compositeness, but those mechanisms require
an explicit comparison with this notebook's target. The next
construction task would specify one candidate lepton operator and
scalar partner, their conserved charges and SUSY realization, then
their current overlaps. No new finite-$Q$ result or flavor fit is
claimed in this source-comparison step.

\section{The 2024 sBootstrap and the SO(32) completion}
\label{sec:so32completion}
The subsequently supplied paper by Rivero~\cite{Rivero2024} makes an
important refinement to the preceding comparison. Sections 5 and 8
explicitly distinguish its proposed scalar composites from ordinary
QCD-scale mesons and diquarks, and envisage a different high-energy
binding mechanism. Thus the distinctive hypothesis is that a subset
of observed SM fermions supplies the constituents of scalar partners;
it is not necessarily an identification of those partners with known
QCD hadrons. This narrows the contrast with new strong dynamics in
Refs.~\cite{ARS2025,AGRS2026}, although the constituent and operator
dictionaries remain different.

\paragraph{What the five-flavor count supplies.}
Let $V=(5,3)$ under $SU(5)_f\times SU(3)_c$.
Its complex dimension is 15, so its underlying real space has dimension
30. The natural orthogonal extension is therefore $SO(30)$.
Its adjoint already contains the channels used in the scalar count:
\begin{align}
435&=224+105+\overline{105}+1,\\
224&=(24,1)+(1,8)+(24,8),\\
105&=(15,\bar3)+(10,6).
\end{align}
Indeed Sections 3.2 and 7.1 of Ref.~\cite{Rivero2024} explicitly
recognize $SO(30)$ as sufficient for the relevant decompositions.

The minimal extension to 32 real labels is one additional complex
color--flavor singlet:
\begin{equation}
16=(5,3)+(1,1),\qquad
32_{\mathbb C}=(5,3)+(\bar5,\bar3)+1+1.
\end{equation}
These are internal representation dimensions, not counts of spinor
components. A sterile-neutrino-like endpoint and its conjugate are a
possible physical interpretation, not a result of the count. Here the
fifteen labels are five quark flavors times three colors, rather than
the fifteen Weyl fermions of one SM family. The familiar right-handed
neutrino completion of a different family embedding therefore does
not automatically apply. Adding the omitted colored top instead
gives 18 complex labels and points to $SO(36)$.

\paragraph{The extra adjoint content.}
Taking the exterior square of $\mathbb R^{30}\oplus\mathbb R^2$
gives the branching under $SO(30)\times SO(2)$,
\begin{equation}
496=(435,1)+(30,2)+(1,1).
\end{equation}
The difference is 61 adjoint components, not just two extra scalar
states. In color--flavor language the added content is
$2(5,3)+2(\bar5,\bar3)+(1,1)$, precisely the mixed channels and
additional singlet visible in Eq.~(1) of Ref.~\cite{Rivero2024}.
A candidate neutral endpoint must account for these mixed colored
channels. Its electric charge also needs an explicit embedding.
All these identities are dimensionless and make no convention
setting physical $\hbar$ to one.

\paragraph{An enlargement is not yet a selection rule.}
The paper suggests replacing color $SU(3)$ by $U(3)$.
However, the unitary action $A\otimes B$ of
$SU(5)\times U(3)$ on $\mathbb C^{15}$ already acts orthogonally
on $\mathbb R^{30}$, with real determinant
$|\det_{\mathbb C}(A\otimes B)|^2=1$.
Its image, up to a finite kernel, lies in $U(15)\subset SO(30)$.
The extra central phase alone consequently does not force 32
real directions. Nor does adding singlets by itself demonstrate
interactions invariant under the full larger algebra.

One possible additional selector is the standard ten-dimensional
$N=1$ supergravity/super-Yang--Mills Green--Schwarz anomaly-cancellation
condition. Its gravitational anomaly requires 496 gauge generators;
within the orthogonal family,
\begin{equation}
\frac{n(n-1)}{2}=496\quad\Longrightarrow\quad n=32.
\end{equation}
Gauge and mixed anomaly factorization are also required; the type-I
string implements a related selection through tadpole cancellation.
This is a conditional higher-dimensional input, not a consequence of
the four-dimensional scalar count. The presence of 496 scalar
components alone does not supply the necessary ten-dimensional
gauge sector, and $E_8\times E_8$ is another familiar string possibility.

The missing ingredient is therefore a physical reason for exactly
one complex neutral sector and for the full symmetry, or a genuine
string-consistency condition selecting it. A neutral constituent is
a useful guess; no particular particle assignment, electroweak
completion or binding dynamics is derived here.

\section{Scalars as internal components of the ten-dimensional gauge field}
\label{sec:gaugefieldscalars}
Section~\ref{sec:so32completion} left two questions open: what removes
the adjoint components outside the selected channels, and why a scalar
multiplet with 496 entries should obey a condition on the
ten-dimensional gauge sector. This section tests one hypothesis that
addresses both: the scalars of Ref.~\cite{Rivero2024} are the components
of the ten-dimensional $SO(32)$ gauge field along compact directions.
All statements are conditional on that identification. Except for
Eq.~\eqref{eq:wilsonmass}, they are dimensionless representation
statements. Details, scripts and outputs are in blackboard~012 of the
research notebook.

\paragraph{Labels.}
Fifteen of the sixteen Cartan directions carry a pair (flavor, color),
with flavor in $\{u,c,d,s,b\}$; the sixteenth, $n$, is colorless.
The roots $\pm e_a\pm e_b$ then split into meson-like roots
$e_{(i,c)}-e_{(j,d)}$, diquark-like roots $e_{(i,c)}+e_{(j,d)}$, which
carry $(15,\bar3)+(10,6)$, and one-label roots $e_{(i,c)}\pm e_n$.
Electric charge is the Cartan generator $Q=\sum_a q_aH_a$ with the quark
charges on the colored labels; it reproduces the charges of
Ref.~\cite{Rivero2024} on the selected channels.

\paragraph{Nine dimensions.}
The nine-dimensional $N=1$ vector multiplet contains one real scalar.
For a string theory on a circle it is the gauge-field component $A_9$,
in the 496. Restricted to $(24,1)+(15,\bar3)+(\overline{15},3)$ it has
$24+2\times45=114$ real components: twelve for charged and twelve for
neutral slepton-like states, and, per color, six complex states of
charge $1/3$, six of charge $2/3$ and three of charge $4/3$. This is the
count of two scalars per Dirac fermion of Ref.~\cite{Rivero2024} with no
additional copies. Each further compact dimension adds one real scalar
per generator. Reference~\cite{Rivero2009} had singled out nine
dimensions as the non-chiral stage. Under the hypothesis the scalar
multiplet is the gauge sector, so the ten-dimensional anomaly condition
of Section~\ref{sec:so32completion} becomes the selector of $SO(32)$
over $SO(30)$.

\paragraph{Label symmetry.}
In the fermionic formulation of the heterotic string the 32 labels are
carried by left-moving worldsheet fermions $\lambda^I$, and the adjoint
states are the bilinears $\lambda^I_{-1/2}\lambda^J_{-1/2}|0\rangle$,
antisymmetric by worldsheet statistics. Space-time spin comes from the
right-movers. The channel $(15,\bar3)=S^2(5)\otimes\Lambda^2(3)$ is
therefore a bilinear of two label-carrying fermions with no constraint
from space-time statistics. The left-moving internal central charge 16
corresponds to sixteen complex labels, $5\times3+1$.

\paragraph{Mass formula.}
With a Wilson line $a$ on a circle of radius $R$, a state on root $p$
has, at the lowest Kaluza--Klein level,
\begin{equation}
M=\frac{\hbar\,|p\cdot a|}{R\,c_{\rm light}},\qquad
M^2\propto(a_a\mp a_b)^2
\label{eq:wilsonmass}
\end{equation}
for meson-like and diquark-like roots respectively. The self-energy
formula $E\propto(q_a+q_b)^2$ of Ref.~\cite{Rivero2024} has this form if
$E$ is read as $M^2$, with the condition $z_1+z_2+z_3=0$ a Wilson line in
the Cartan subalgebra of the $SU(3)$ acting on $(d,s,b)$. The second
condition of that formula is an additional input. A label reflection
$e_a\to-e_a$ is a Weyl reflection and exchanges meson-like and
diquark-like sectors. In nine dimensions a scalar on a root with
$p\cdot a\neq0$ is absorbed by the massive vector.

\paragraph{Orbifold projections.}
For $T^6/\mathbb Z_N$ orbifolds preserving four-dimensional $N=1$, with
heterotic consistency conditions, no Wilson lines and color-blind shifts
$V=(w_u^3,w_c^3,w_d^3,w_s^3,w_b^3,w_n)$, a root gives a gauge boson if
$p\cdot V\in\mathbb Z$, one chiral multiplet from $A_{z_k}$ if
$p\cdot V=v_k$ mod 1, and is projected out otherwise. An exact
enumeration gives the following.
\begin{itemize}
\item $\mathbb Z_3$ reproduces the series $U(3k)\times SO(32-6k)$.
Its member with gauge group $U(15)\times SO(2)$ has untwisted matter
$3\,[\overline{105}+15_{+1}+15_{-1}]$ and twisted matter
$27\,[15_{-1/2}+3\cdot1]$: three copies of the diquark-like multiplet,
one per complex plane.
\item Single copies occur. With the split $(u,c\,|\,d,s,b)$, shifts in
$\mathbb Z_8$-I and $\mathbb Z_{12}$-I give $U(6)\times U(9)\times U(1)$
with every diquark-like label as one chiral multiplet and the charged
meson-like labels projected out.
\item The $\mathbb Z_7$ shift $V=(1^6,2^9,0)/7$ gives
$U(6)\times U(9)\times SO(2)$ with every charged label exactly once
(Table~\ref{tab:z7}); the twelve neutral meson-like labels are gauge
bosons of the flavor subgroup.
\item No admissible shift isolates the selected channels: $(10,6)$
accompanies $(15,\bar3)$ and $(24,8)$ accompanies $(24,1)$. Among the
96 shifts that make all fifteen diquark-like labels chiral with one
orientation, none keeps a meson-like label as matter.
\end{itemize}
Cubic non-abelian anomalies cancel between untwisted and twisted sectors
in the five non-trivial $\mathbb Z_3$ models and in the $\mathbb Z_7$
model, which fixes the orientation convention. Twisted sectors of the
non-prime orbifolds were not computed.

\begin{table}[ht]
\centering
\caption{Massless charged spectrum of the $\mathbb Z_7$ orbifold with
$V=(1^6,2^9,0)/7$. Untwisted states are internal components of the
ten-dimensional gauge field.}
\label{tab:z7}
\begin{tabular}{@{}llll@{}}
\toprule
Sector & Copies & $U(6)\times U(9)$, $SO(2)$ charge & Labels \\
\midrule
untwisted, plane 1 & 1 & $(\bar6,9)$ & charged mesons $d\bar u$ \\
untwisted, plane 1 & $1+1$ & $(6,1)_{\pm1}$ & $u\,n$, $u\,\bar n$ \\
untwisted, plane 2 & 1 & $(15,1)$ & $uu$, charge $4/3$ \\
untwisted, plane 2 & $1+1$ & $(1,9)_{\pm1}$ & $d\,n$, $d\,\bar n$ \\
untwisted, plane 3 & 1 & $(1,36)$ & $dd$, charge $-2/3$ \\
untwisted, plane 3 & 1 & $(\bar6,\bar9)$ & $\bar u\bar d$, charge $-1/3$ \\
twisted $\theta$ & 21 & singlets & \\
twisted $\theta^2$ & 7, 7, 21 & $(6,1)$, $(1,\bar9)$, singlets & \\
twisted $\theta^4$ & 7, 7 & $(6,1)$, singlets & \\
\bottomrule
\end{tabular}
\end{table}

\paragraph{The sixteenth label.}
The 61 generators that distinguish $SO(32)$ from $SO(30)$ are the
states with exactly one colored label and one singlet. With
$\mathfrak{so}(32)=\mathfrak k\oplus\mathfrak p$,
$\mathfrak k=\mathfrak{so}(30)\oplus\mathfrak{so}(2)$ and
$\mathfrak p=(30,2)$, one has
$[\mathfrak p,\mathfrak p]\subset\mathfrak k$: two one-label states
compose into a meson-like or diquark-like label. The charge formula of
Ref.~\cite{Rivero2024} gives the new label charge zero. If instead $Q$
is required to be traceless on the sixteen complex labels, the colored
labels sum to $3(2/3+2/3-1/3-1/3-1/3)=1$ and the new label has charge
$-1$. The one-label states are then leptoquark-like. Three of them,
$\bar d\ell^+$, $\bar s\ell^+$, $\bar b\ell^+$, have color $\bar3$,
charge $4/3$ and $B-L=2/3$, the quantum numbers of the three $uu$-type
states, and would complete that sector to six. A colored top would
instead require eighteen labels. The tracelessness requirement is an
assumption, and the same label adds two states of charge $1/3$, three
of charge $2/3$ and two of charge $5/3$ in color $\bar3$.

\paragraph{What is missing.}
The gauge groups found are $U(15)$ or $U(6)\times U(9)$; the product
structure $15=5\times3$ is a choice of basis, and no mechanism here
leaves $SU(3)$ color gauged with flavor global. Ordinary supersymmetry
pairs each scalar with a fermion of the same labels, whereas the
sBootstrap pairs a two-label scalar with a one-label fermion. Abelian
anomalies, Wilson lines and couplings were not analyzed.

\section{Color from an 't Hooft twist}
\label{sec:colortwist}
Section~\ref{sec:gaugefieldscalars} ended with gauge groups $U(15)$ or
$U(6)\times U(9)$, in which $15=5\times3$ is a choice of basis. This
section examines one mechanism that makes the product structure
physical: a flux background of the gauge field on a two-torus, the
setting of Ref.~\cite{CDS} at rational noncommutativity. Statements
about twisted bundles and index numbers are standard and are rederived
in blackboard~013; the conclusions for the sBootstrap are conditional.

\paragraph{Twist.}
A $U(15)$ bundle of first Chern number $m=3m'$, with $m'$ prime to 5,
has transition functions $U\otimes1_3$ and $V^{m'}\otimes1_3$ on adjoint
fields, with $U,V$ the $5\times5$ clock and shift matrices. Their
commutant is $1_5\otimes U(3)$: color is the unbroken group and the
five-fold index is a wrapping number.

\paragraph{Mesons on a noncommutative torus.}
Expanding an adjoint field in $U^aV^b\otimes\phi_{ab}(x)$ gives the
$3\times3$ fields $\phi_{ab}$ fractional momenta $b/5$ and $-m'a/5$ in
units of $2\pi\hbar/L_i$. The single-valued generators
$Z_1=Ve^{2\pi ix_1/5L_1}$ and $Z_2=U^{-s}e^{2\pi ix_2/5L_2}$, with
$sm'=1$ mod 5, satisfy $Z_1Z_2=e^{2\pi is/5}Z_2Z_1$. The meson-like
fields, read in Ref.~\cite{Rivero2024} as $(24+1)\otimes(8+1)$, are
therefore $U(3)$ fields on the noncommutative torus with $\theta=s/5$:
the number of colors is the rank and the number of flavors is the
denominator of $\theta$. Flavor symmetry is reduced to
$\mathbb Z_5\times\mathbb Z_5$, and the lowest state of class $(a,b)$
has
\begin{equation}
M=\frac{2\pi\hbar}{c_{\rm light}}
\sqrt{\Bigl(\frac{\tilde b}{5L_1}\Bigr)^2+\Bigl(\frac{\tilde a}{5L_2}\Bigr)^2},
\end{equation}
with $\tilde a,\tilde b\in\{-2,\dots,2\}$ the representatives of $-m'a$
and $b$. On a square torus the 24 non-singlet classes split as
$4,4,4,8,4$ at $M^2\propto1,2,4,5,8$; only the $U(3)$ gauge multiplet
keeps a zero mode.

\paragraph{Diquark-like zero modes.}
Since $\Lambda^2(5\otimes3)=S^2(5)\otimes\Lambda^2(3)+
\Lambda^2(5)\otimes S^2(3)$, the numbers of color-antitriplet and
color-sextet zero modes on the torus are $c_1(S^2E_5)=6m'$ and
$c_1(\Lambda^2E_5)=4m'$, in place of the flavor multiplicities 15 and
10. For wrapping number $n$ they are $(n+1)m'$ and $(n-1)m'$.

\paragraph{Charge.}
For any background of the form $f\otimes1_3$ the unbroken group is the
unitary group of $S'\otimes M_3$, with $S'$ the commutant of the flavor
parts. A single color group needs $S'=\mathbb C$, which leaves no
flavor-dependent charge. A charge with eigenvalues $2/3$ and $-1/3$
needs $S'\supseteq\mathbb C\oplus\mathbb C$ and gives
$U(3)_u\times U(3)_d$. The corresponding construction has two stacks
with wrapping numbers 2 and 3, fluxes $m_u$ and $m_d$, the label count
$3\cdot2+3\cdot3+1=16$, and $Q=\tfrac23Q_u-\tfrac13Q_d+q_nQ_n$. Its
zero modes on the flux torus are $3|m_u|$ antitriplets and $|m_u|$
sextets in the $uu$ sector, $4|m_d|$ and $2|m_d|$ in the $dd$ sector,
$|3m_u+2m_d|$ $ud$ diquarks and $|3m_u-2m_d|$ charged mesons. For
$m_u=1$ the $uu$ sector has the three antitriplets of
Ref.~\cite{Rivero2024}; the counts six and six of the other sectors
would need $m_d=3/2$. A scan of factorized fluxes on $T^6$ with flux
numbers between $-4$ and $4$ finds no assignment with the ratio
$1:2:2$ and none that cancels the tadpoles with these three stacks.

\paragraph{Assessment.}
The flux mechanism singles out color by turning flavor into wrapping.
The multiplicities that Ref.~\cite{Rivero2024} counts as dimensions of
flavor representations become index numbers and momentum classes with
different values, and a flavor-dependent charge requires two color
groups. The counts of Ref.~\cite{Rivero2024} are recovered when the flux
is absent, where color is not singled out. A background outside the
flavor directions, or a mechanism other than symmetry breaking, would
be needed to obtain both.

\section{Three-label states: baryon labels in \texorpdfstring{$E_8$}{E8} and in the spinor class}
\label{sec:threelabel}
The adjoint of $SO(32)$ organizes two-label states. This section records
what $E_8$ holds instead, and where three-label states sit on the
$SO(32)$ side. Details and the exact-arithmetic script are in
blackboard~014. The representation content is computed; its reading in
terms of flavor, color and charge is conditional, and statements about
known string models and dualities are quoted from memory.

\paragraph{$E_8$ in $SU(9)$ form.}
The roots of $E_8$ are $e_i-e_j$ and
$\pm(e_i+e_j+e_k-\tfrac13\sum e)$ in $\mathbb R^9$, so that
$248=80+84+\overline{84}$ with $84=\Lambda^3(9)$. Reading the nine
labels as three flavors times three colors,
\begin{equation}
84=(10,1)+(8,8)+(1,10),\qquad 80=(8,1)+(1,8)+(8,8).
\end{equation}
The $(10,1)=S^3(3_f)\otimes\Lambda^3(3_c)$ is the label content of the
baryon decuplet. For $(u,d,s)$ the charge
$Q=\mathrm{diag}(2/3,-1/3,-1/3)\otimes1_3$ is traceless and is a
generator of $E_8$; the color singlets of the 84 then have charges
$2,1,0,-1$ with multiplicities $1,2,3,4$, and those of the 80 have
charges $-1,0,1$ with multiplicities $2,4,2$. The decomposition is a
$\mathbb Z_3$ grading: two elements of the 84 bracket into the
$\overline{84}$, so label number is conserved modulo 9. There is no
$\Lambda^2(9)$, hence no diquark-like state, and no baryon octet, which
would need spin labels and $SU(18)$.

\paragraph{A string model.}
Consider the $\mathbb Z_3$ orbifold of the $E_8\times E_8$ heterotic
string with shift
\begin{equation}
V=\tfrac13\,(1,1,1,1,2,0,0,0\,;\,2,0,0,0,0,0,0,0).
\end{equation}
Its gauge group is $SU(9)\times SO(14)\times U(1)$, its untwisted matter
is $3\,[(84,1)+(1,14)+(1,64)]$ and its twisted matter is
$27\,(\bar9,1)$. The $SU(9)$ cubic anomaly cancels as $3\cdot9-27$. In the label reading the
untwisted matter carries baryon-like labels and the twisted matter
anti-label states, where the $SO(32)$ model of
Section~\ref{sec:gaugefieldscalars} has diquark-like and one-label
states.

\paragraph{The spinor class of $SO(32)$.}
The lattice of $\mathrm{Spin}(32)/\mathbb Z_2$ contains the weights
$(\pm\tfrac12)^{16}$ with an even number of minus signs, that is,
subsets $S$ of the sixteen labels with $|S|$ even. These states are
massive at the string scale. Among them the color singlets have
$0,3,6,9,12,15$ colored labels, with multiplicities
$1,35,175,175,35,1$, and the sixteenth label is present exactly when the
baryon number $B$ is odd. Three consequences follow. Baryon labels exist
in the $SO(32)$ theory, as the $35=S^3(5)$, together with a tower of
multi-baryon label states. The lattice condition reads $3B+L$ even if
the sixteenth label carries lepton number one. And since a subset has
charge $\sum_{a\in S}q_a-\tfrac12\sum_{\rm all}q_a$, with the colored
labels summing to one, the color singlets have integer charge exactly
when the sixteenth label has charge $-1$. This supplies the trace
condition that Section~\ref{sec:gaugefieldscalars} had to assume.

\paragraph{Remarks.}
The decomposition $80+84+\overline{84}$ also organizes the U-duality
algebra $E_{8(8)}\supset SL(9,\mathbb R)$ of M-theory on an eight-torus,
with the 84 holding the three-form components, to which the membrane
couples, and the $\overline{84}$ those of the dual six-form of the
fivebrane; Ref.~\cite{Rivero2009} had matched that 84 with the light
fermion degrees of freedom. This is the same complex algebra in a
different real form and role, and no connection between the two uses is
derived here. An earlier suggestion of this notebook session, that the
$U(9)$ of nine $SO(32)$ labels could enhance to $E_8$ in nine
dimensions, is withdrawn: the duality between the two heterotic frames
does not carry labels to labels, and in the $SO(32)$ frame the
three-label states are those of the spinor class.

\section{Nine dimensions: winding labels and the exceptional groups}
\label{sec:ninedim}
On a circle the two heterotic strings are one theory, so the label
descriptions of Sections~\ref{sec:gaugefieldscalars}
and~\ref{sec:threelabel} can be compared. This section records the
comparison. Details and the exact-arithmetic script are in
blackboard~015. The lattice statements are computed. The diagram of the
lattice and the list of points of maximal symmetry are
known~\cite{FFGNP}. The reading in terms of flavor, color and charge is
conditional, and remarks on the type~I$'$ description are quoted from
memory.

\paragraph{Setting.}
Charges form the even self-dual lattice of signature $(17,1)$. In
$SO(32)$ variables a vector is $v=(\pi;n,w)$, with $\pi$ the sixteen
label charges, $n$ the momentum and $w$ the winding number, and
$v^2=\pi^2+2nw$. A vacuum is a Wilson line $A$ and a radius $R$,
combined in $t=(-A;-(A^2/2+r),1)$ with $r=R^2/\alpha'$. A vector with
$v^2=2$ is a state with
\begin{equation}
\alpha' M^2=\frac{(v\cdot t)^2}{r}
\label{eq:ninedmass}
\end{equation}
in units $\hbar=c=1$ and with physical label charges $\pi+wA$; the
massless ones are the charged gauge bosons. A set of nineteen simple
roots is $g_1=(0;-1,-1)$, $a_0=(-e_1-e_2;1,0)$,
$a_i=(e_i-e_{i+1};0,0)$ for $i=1,\dots,15$,
$a_{16}=(e_{15}+e_{16};0,0)$ and $g_2=((-\tfrac12)^{16};1,-1)$. Two of
them carry winding. Deleting two nodes of their diagram gives 44 points
of maximal symmetry up to its left--right symmetry, and a direct
enumeration of roots agrees with the diagram at each.

\paragraph{Exceptional groups.}
With $m$ labels at Wilson-line value $1/2$, the others at $0$ and
$r=(8-m)/8$, the gauge group is $E_{m+1}\times SO(32-2m)$, where
$E_4=SU(5)$ and $E_5=SO(10)$. The exceptional factor consists of the
$SO(2m)$ of the $m$ labels at zero winding, its spinor at winding
$\pm1$ and, for $m=6$ and $7$, a singlet or the vector at winding
$\pm2$. For $m=7$,
\begin{equation}
248=14+64+92+64+14\qquad(w=-2,\dots,2),
\end{equation}
and the winding-one states are the subsets of the seven labels of even
size. If the seven labels are $(u,c)$ times color and the sixteenth
label, three colored labels come with the sixteenth. The color singlets
at winding one are then the empty set, the four flavor-symmetric states
$uuu$, $uuc$, $ucc$, $ccc$ with the sixteenth label, and the state with
all six colored labels.

\paragraph{The nine labels of $E_8$.}
An $SU(9)$ inside this $E_8$ has seven $SO(32)$ labels and two labels
that differ from them by winding: the nine labels carry winding
$(\tfrac13)^7,-\tfrac23,-\tfrac53$. The three-label multiplet splits as
$84=35+21+21+7$ at winding $1,0,-1,-2$. The 21 at zero winding are
$SO(32)$ two-label states, among them the fifteen diquark-like states of
the six colored labels. A diquark-like state of $SO(32)$ is thus a
three-label state of $E_8$ whose third label is a winding label. Other
duality frames grade $E_8$ differently: among 6000 frames generated by
reflections its zero-winding part is $SO(14)$, $SU(8)$, $SO(12)$,
$SU(7)\times SU(2)$, $SU(7)$ or $SO(10)\times SU(3)$. With $SU(8)$,
eight of the nine labels are $SO(32)$ labels of that frame. This
replaces by a computation the argument at the end of
Section~\ref{sec:threelabel}.

\paragraph{Color.}
In any frame of $SO(32)$ type the roots of $E_8$ at zero winding embed
in the roots of $D_{16}$, which has no subsystem of type $E$. For nine
labels read as $3_f\otimes3_c$, every color-blind direction is orthogonal
to the roots of an $E_6$. Hence the winding number never commutes with
the color group of Section~\ref{sec:threelabel}, and the nine labels
are never all $SO(32)$ labels: the decuplet reading of $E_8$ has no
expression in $SO(32)$ labels. The color group of
Section~\ref{sec:gaugefieldscalars} behaves differently. At the
$E_8\times SO(18)$ point it sits at zero winding and acts on the nine
labels as $3+3+1+1+1$. Its singlets in the 248 number 22 and form
$G_2\times SU(3)$. The $SU(3)$ acts on the three colorless labels, and
the long roots of $G_2$ are $uuu$, $ccc$, the state of the three
colorless labels and their conjugates.

\paragraph{Vacua that respect color and charge.}
Of the 44 points, 24 keep the color group of
Section~\ref{sec:gaugefieldscalars} at zero winding. For 11 of them the
Wilson line is constant on the six $(u,c)$ labels and on the nine
$(d,s,b)$ labels (Table~\ref{tab:ninedim}). In that class exceptional
symmetry arises only as $E_7$ or $E_8$ and only on the up-type block,
with the sixteenth label deciding between the two.

\begin{table}[ht]
\centering
\caption{Points of maximal gauge symmetry of the heterotic string on a
circle whose Wilson line depends only on electric charge. The last
column gives the Wilson-line value of each block of labels, folded to
$[0,\tfrac12]$; $n$ is the sixteenth label.}
\label{tab:ninedim}
\begin{tabular}{@{}llll@{}}
\toprule
Blocks of labels & Gauge group & $R^2/\alpha'$ & Wilson line \\
\midrule
all sixteen & $SO(32)\times SU(2)$ & $1$ & $0$ \\
 & $SO(34)$ & $1/2$ & $0$, sign on the spinor class \\
 & $SU(16)\times SU(2)\times SU(2)$ & $1/2$ & $1/4$ \\
 & $SU(17)\times SU(2)$ & $17/49$ & $3/14$ \\
 & $SU(18)$ & $1/4$ & $1/4$ \\
$(u,c,n\,|\,d,s,b)$ & $E_8\times SO(18)$ & $1/8$ & $1/2\,|\,0$ \\
 & $E_8\times SU(9)\times SU(2)$ & $1/9$ & $1/2\,|\,1/18$ \\
 & $E_8\times SU(10)$ & $5/49$ & $1/2\,|\,1/14$ \\
$(u,c\,|\,d,s,b,n)$ & $E_7\times SO(20)$ & $1/4$ & $1/2\,|\,0$ \\
 & $E_7\times SU(10)\times SU(2)$ & $1/5$ & $1/2\,|\,1/10$ \\
 & $E_7\times SU(11)$ & $11/64$ & $1/2\,|\,1/8$ \\
\bottomrule
\end{tabular}
\end{table}

\paragraph{Massive three-label states.}
At the $E_8\times SO(18)$ point the lightest massive states have
$\alpha'M^2=1/2$ and form the spinor of $SO(18)$ at winding $\pm1$.
They are vectors of the spinor class of Section~\ref{sec:threelabel}
with one unit of winding, and they contain $\Lambda^3(9)$ of the nine
$(d,s,b)$ labels, with the ten decuplet labels as color singlets. At the
$SU(18)$ point the eighteen labels are the sixteen $SO(32)$ labels and
two winding labels. Every state there has label number divisible by
three, so the group is $SU(18)/\mathbb Z_3$, and the lightest massive
states, at $\alpha'M^2=1$, are the 816 states of $\Lambda^3(18)$ and
their conjugates. Of these, 120 have zero winding and are the
diquark-like states of $SO(32)$.

\paragraph{Strings that are points in the noncompact directions.}
Two mechanisms of this setting leave a string without extension in the
noncompact directions; blackboard~016 records them. First, a Wilson
line proportional to electric charge, $A=\tfrac32Q$, has holonomy
$(-1)^{3Q}$. A two-label state is lifted exactly when $3Q$ is odd, with
$\alpha'M^2=1/(4r)$. Of the 24 color-singlet meson-like labels the 12
neutral ones stay massless and the 12 of charge $\pm1$ are lifted. In
the type~I$'$ description a lifted state is an open string stretched
across the compact interval. For charge $-1$ of the sixteenth label this
vacuum is the $E_7\times SO(20)$ point at $r=1/4$. Second, the winding
states are D-particles in that description, and an open string between
a D-particle and a D8-brane has mixed boundary conditions in eight
directions: it has no position of its own and its only massless state
is fermionic~\cite{BGS,BGL}. At the $SU(18)$ point the 32 massless
states of winding one are single labels attached to a winding label:
fifteen quark-like color triplets and one state with the sixteenth
label and charge $-1$, with their conjugates, for the traceless charge
$\mathrm{diag}(0;q_1,\dots,q_{16};0)$. Such a state keeps the
zero-point motion of its vibration modes, a mean square displacement
$\alpha'\sum_{n\le N}1/n$ per noncompact direction, set by the slope of
the string alone and independent of its internal length. For
the fundamental string of this setting that spread is far below any
bound, so these states are point-like. The hadronic length of the
bosonic partners must then enter by another route: this is the question
of scales of the earlier sections seen from the other side.

\paragraph{What is missing.}
These vacua have sixteen supercharges and only adjoint matter. Nothing
here selects a vacuum, produces chirality or fixes the charge of
winding states, which needs one number beyond the label charges.


\section*{Reproducibility and review scope}
The radius derivation is analytic; no numerical physics script was required. A research agent derived the leading anharmonic radius and the parent independently checked only that new mathematics. The parent analyzed the conceptual $B-L$ ledger and linear-multiplet constraint. Earlier source reading and document builds were delegated to Sol. The subsequent representation-count clarification and 1983 source comparison add no finite-$Q$ calculation. Version 0.6 added the parent's reading of supplied 2025/2026 LaTeX sources and the scoped comparison in Section~\ref{sec:partialcomparison}. Version 0.7 corrects the source interpretation using the supplied 2024 paper and adds the dimensionless representation analysis in Section~\ref{sec:so32completion}. Sol extracted that PDF and validates the document; the parent checked consequential passages on original page images. Version 0.8 adds Section~\ref{sec:gaugefieldscalars}. It was written on 10 November 2026 by Claude (model Fable 5.2) in a separate session, after reading Refs.~\cite{Rivero2024,Rivero2005,Rivero2009}, a 2007 preprint of the same program and Section~\ref{sec:so32completion}. Its enumeration of orbifold shifts and its twisted spectra use exact-arithmetic scripts in \path{scripts/so32_gauge_scalars/}, written and run by the same model, with outputs in \path{calculations/so32_gauge_scalars/}. The independent consistency criteria are the reproduction of the $\mathbb Z_3$ series of gauge groups and the cancellation of cubic anomalies between sectors computed by different rules. No second model reviewed that section, and it was built with \texttt{pdflatex} because \texttt{latexmk} was unavailable. Version 0.9 adds Section~\ref{sec:colortwist}, written the same day by the same model after the user pointed to Ref.~\cite{CDS}; only the abstract of that reference was retrieved, and the standard results it relies on are quoted from memory and rederived in blackboard~013. Its sanity checks, tables and bounded flux scan come from \path{scripts/so32_gauge_scalars/flux_twist.py}. Version 0.10 adds Section~\ref{sec:threelabel}, same day, model and session, after the user asked for the observation on baryon labels in $E_8$ to be worked out; its decompositions, singlet counts and orbifold spectrum come from \path{scripts/e8_baryons/e8_baryons.py}, with weight-multiset comparison and anomaly cancellation as consistency checks. Version 0.11 adds Section~\ref{sec:ninedim}, same day, model and session, written at the session's own initiative as a follow-up of an open point of blackboard~014 and not separately requested by the user; its lattice enumerations come from \path{scripts/nine_dim_frames/nine_dim_frames.py}, with the agreement between the diagram and a direct enumeration of roots at 80 points, and the dimensions of the multiplets found, as consistency checks. Version 0.12 adds one paragraph to that section, after the user asked whether a compactification can force a fermionic string to be a point in the noncompact directions; its two lattice checks come from \path{scripts/nine_dim_frames/charge_wilson_line.py}, and its D-brane statements are quoted from memory. Version 0.13 rewrites the closing sentences of that paragraph, on the size of such states, after an objection of the user. Earlier recoil, energy and current calculations are unchanged. These checks do not establish experimental correctness, novelty or a completed supersymmetric model.

The living source and PDF are in
\begin{quote}\small
\path{discussions/2026-10-11-quark-antiquark-open-string/manuscript/}
\end{quote}
The research notebook records the model assumptions, source exposure, prior results and next task. The historical motivations include Schwarz's dual quark--gluon model~\cite{Schwarz}, with Seiberg duality~\cite{Seiberg} retained as a possible later connection. Neither reference proves the contact interaction or its intended string interpretation.

\begin{thebibliography}{99}
\bibitem{Rivero2024} A. Rivero,
\emph{An interpretation of scalars in SO(32)},
Eur. Phys. J. C \textbf{84}, 1058 (2024),
\href{https://doi.org/10.1140/epjc/s10052-024-13368-3}{doi:10.1140/epjc/s10052-024-13368-3}.
Supplied original PDF; representation and discussion sections reviewed.
\bibitem{CDS} A. Connes, M. R. Douglas and A. Schwarz,
\emph{Noncommutative geometry and matrix theory: compactification on tori},
JHEP \textbf{02} (1998) 003,
\href{https://arxiv.org/abs/hep-th/9711162}{arXiv:hep-th/9711162}.
Metadata and abstract retrieved; the full text was not re-read for this version.
\bibitem{FFGNP} A. Font, B. Fraiman, M. Gra\~na, C. A. N\'u\~nez and H. Parra De Freitas,
\emph{Exploring the landscape of heterotic strings on $T^d$},
JHEP \textbf{10} (2020) 194,
\href{https://arxiv.org/abs/2007.10358}{arXiv:2007.10358}.
Located by a search; the text was not re-read for this version.
\bibitem{BGS} C. P. Bachas, M. B. Green and A. Schwimmer,
\emph{(8,0) quantum mechanics and symmetry enhancement in type I$'$ superstrings},
JHEP \textbf{01} (1998) 006,
\href{https://arxiv.org/abs/hep-th/9712086}{arXiv:hep-th/9712086}.
Abstract retrieved by a search; the statement cited is quoted from memory.
\bibitem{BGL} O. Bergman, M. R. Gaberdiel and G. Lifschytz,
\emph{String creation and heterotic--type I$'$ duality},
Nucl. Phys. B \textbf{524} (1998) 524,
\href{https://arxiv.org/abs/hep-th/9711098}{arXiv:hep-th/9711098}.
Abstract retrieved by a search; the statement cited is quoted from memory.
\bibitem{ARS2025} K. Agashe, L. Ricci and R. Sundrum,
\emph{Partial Compositeness: from Anarchy to Symmetry},
\href{https://arxiv.org/abs/2507.05332v1}{arXiv:2507.05332v1} (2025).
Supplied LaTeX source, especially Sections 2--4.
\bibitem{AGRS2026} K. Agashe, G. F. Giudice, R. Rattazzi and R. Sundrum,
\emph{The Super-Composite Higgs},
\href{https://arxiv.org/abs/2609.26413v1}{arXiv:2609.26413v1} (2026).
Supplied LaTeX source; construction, hidden-SQCD toy and summary reviewed.
\bibitem{Schwarz} J. H. Schwarz, \emph{Dual quark-gluon model of hadrons},
Phys. Lett. B \textbf{37}, 315--319 (1971),
\href{https://doi.org/10.1016/0370-2693(71)90028-1}{doi:10.1016/0370-2693(71)90028-1}.
Metadata and abstract verified previously; full primary text not yet reviewed.
\bibitem{Seiberg} N. Seiberg, \emph{Electric-magnetic duality in supersymmetric non-Abelian gauge theories},
\href{https://arxiv.org/abs/hep-th/9411149}{arXiv:hep-th/9411149}.
Relevant sections were read in the preceding notebook work; no new operator matching is claimed here.
\bibitem{Rivero2005} A. Rivero, \emph{Supersymmetry with composite bosons},
\href{https://arxiv.org/abs/hep-ph/0512065}{arXiv:hep-ph/0512065} (2005).
The primary PDF was read; its degree count and proposed lepton-sector extension are prior art, not a completed dynamics.
\bibitem{Rivero2009} A. Rivero, \emph{Unbroken supersymmetry, without new particles},
\href{https://arxiv.org/abs/0910.4793}{arXiv:0910.4793} (2009).
The primary PDF was read; its six charged pseudoscalar states and pointlike-lepton objection are relevant here.
\bibitem{GMY1983} O. W. Greenberg, R. N. Mohapatra and M. Yasue,
\emph{Determination of the Number of Generations of Quarks and Leptons from Flavor-Color Symmetry},
Phys. Rev. Lett. \textbf{51}, 1737--1740 (1983),
\href{https://doi.org/10.1103/PhysRevLett.51.1737}{doi:10.1103/PhysRevLett.51.1737}.
The supplied four-page original was read; its conditional anomaly result concerns a distinct supersymmetric preon theory.
\bibitem{VolkovAkulov} D. V. Volkov and V. P. Akulov,
\emph{Is the Neutrino a Goldstone Particle?}, Phys. Lett. B \textbf{46}, 109--110 (1973),
\href{https://doi.org/10.1016/0370-2693(73)90490-5}{doi:10.1016/0370-2693(73)90490-5}.
The supplied two-page original was read; it does not make the pion--muon identification proposed here.
\end{thebibliography}
\end{document}
